\chapter{Discusión}
\label{ch:discusion}

% Integrar los resultados en el marco teórico presentado en la
% Introducción. Comparar con trabajos previos. Limitaciones del trabajo.

En la presente tesis se diseñaron, extrajeron y evaluaron variables a nivel de fijación a partir de los mapas generados por el modelo nnELM, y se estudió su relación con la señal EEG registrada durante una tarea de búsqueda visual híbrida. Los resultados muestran que un subconjunto de las variables latentes del modelo presenta correlatos estadísticamente significativos en la actividad cerebral. En escala original, las variables con clusters significativos según el criterio TFCE fueron \feat{posterior}, \feat{saliencia\_media}, \feat{saliencia\_pixel} y \feat{similitud\_pixel}. Al incorporar transformaciones logarítmicas, este conjunto se amplió: \feat{similitud\_media}, \feat{ganancia\_informacion} y \feat{gap\_entropia} revelaron efectos que no eran detectables en su escala original, lo que indica que la relación entre estas variables y la señal EEG es esencialmente no lineal. En conjunto, se encontraron correlatos significativos en variables pertenecientes a cuatro de los cinco grupos conceptuales definidos: \textbf{información bottom-up}, \textbf{similitud con el target}, \textbf{estado interno del proceso bayesiano} y \textbf{calidad de la decisión}.

Para las variables que no obtuvieron clusters significativos en ninguna de sus versiones evaluadas se pueden realizar distintas interpretaciones. Una de las posibilidades es que estas variables capturen procesos cognitivos cuya expresión neural no es lineal ni capturada adecuadamente por el modelo empleado. En ese caso, métodos no lineales de análisis podrían revelar efectos que no se detectaron en esta oportunidad. 

El objetivo de construir un modelo combinado en el Experimento 3 era evaluar si la contribución individual observada en los Experimentos 1 y 2 se sostenía cuando las variables debían competir por explicar la misma señal EEG; sin embargo, ninguna de las combinaciones evaluadas (Experimentos 3a, 3b y 3c) logró superar en $\Delta$AIC al modelo que incorporaba únicamente \feat{saliencia\_pixel\_log} de forma individual. Más aún, dos de las cuatro variables originalmente seleccionadas por su desempeño individual y su significancia en TFCE (\feat{posterior} y, luego, \feat{similitud\_pixel\_log}) perdieron su cluster espacio-temporal significativo al ser incorporadas junto con las demás, quedando finalmente un modelo reducido a \feat{saliencia\_pixel\_log} y \feat{gap\_entropia\_log} (Experimento 3c).

Quedan abiertas algunas líneas para profundizar este hallazgo en trabajos futuros. En primer lugar, un análisis de partición de varianza permitiría cuantificar qué proporción de la mejora de $\Delta$AIC de cada variable es exclusiva y cuál es compartida con las demás. En segundo lugar, evaluar el orden de incorporación de las variables mediante un esquema \textit{forward stepwise} revelaría si el resultado depende de qué variable ingresa primero al modelo, dado que \feat{saliencia\_pixel\_log} podría estar capturando tempranamente varianza que las otras variables explicarían igualmente bien si se incorporaran solas. 

Cabe destacar, además, que el criterio empleado para construir el subconjunto de variables del Experimento 3 es una entre varias metodologías posibles para construir modelos de mayor complejidad, y estuvo guiado principalmente por los resultados individuales de los Experimentos 1 y 2. Otras estrategias podrían explorar el espacio de combinaciones de manera menos dependiente de esta secuencia de filtros. Asimismo, el presente trabajo no exploró términos de interacción entre variables, los cuales podrían capturar efectos conjuntos no aditivos que un modelo lineal con efectos principales no logra representar. Explorar estas alternativas permitiría no solo obtener un modelo combinado potencialmente más informativo, sino también entender mejor si el resultado observado en el Experimento 3 es un artefacto del criterio de selección utilizado o un reflejo genuino de la estructura de la señal EEG explicada por estas variables.

Un hallazgo relevante que abre una línea de trabajo futura es la presencia de activación significativa en el período previo al onset de la fijación en las variables de saliencia. Esta activación  sugiere que la información de saliencia bottom-up podría estar siendo procesada en relación con la preparación del movimiento ocular más que con el procesamiento del estímulo fijo. En trabajos futuros podría abordarse este fenómeno modelando estas variables en relación con el onset de la sacada en lugar del onset de la fijación, siguiendo la lógica del predictor de intercept de sacada ya incluido en el modelo base.

Una limitación central de este trabajo es que las variables utilizadas como predictores no fueron extraídas directamente del comportamiento del participante sino de las variables latentes de un modelo computacional. El nnELM es una aproximación al proceso de búsqueda humana: si bien replica métricas de comportamiento con alto rendimiento \citep{ruarte2025}, sus mapas internos son construcciones del modelo que pueden diferir del proceso neural real. En consecuencia, la ausencia de correlatos para una variable determinada no implica necesariamente que el proceso cognitivo asociado no tenga representación neural, sino que el modelo y el cerebro pueden estar resolviendo el mismo problema por caminos distintos. Esta limitación subraya la importancia de interpretar los resultados como evidencia de la correspondencia entre el modelo y la actividad cerebral, y no como evidencia directa sobre los mecanismos neurales de la búsqueda.

Por último, todas las variables extraídas en este trabajo se calculan sobre el posterior del target ganador seleccionado por el criterio MinEntropy, sin capturar explícitamente la dinámica de competencia entre los targets sostenidos simultáneamente en memoria durante la búsqueda híbrida. Diseñar variables específicas para esta dimensión constituye una línea de trabajo futura, que permitiría complementar el efecto ya observado para el MSS con una caracterización más fina del proceso de búsqueda híbrida.
